\documentclass{ieeeaccess}
\usepackage{cite}
\usepackage{amsmath,amssymb,amsfonts}
\usepackage{algorithmic}
\usepackage{graphicx}
\usepackage{textcomp}
\usepackage{booktabs}
\usepackage{multirow}
\graphicspath{{../fig/}}
\def\BibTeX{{\rm B\kern-.05em{\sc i\kern-.025em b}\kern-.08em
    T\kern-.1667em\lower.7ex\hbox{E}\kern-.125emX}}

% ---- numbers taken from the tool reports and measurements (updated in one place) ----
\newcommand{\resLUT}{1191}
\newcommand{\resFF}{1824}
\newcommand{\resDSP}{1}
\newcommand{\resBRAM}{2}
\newcommand{\resIOB}{21}
\newcommand{\resPeriod}{7.281}
\newcommand{\resWNS}{+2.719}
\newcommand{\resFmax}{137}

\begin{document}
\history{Date of publication xxxx 00, 0000, date of current version xxxx 00, 0000.}
\doi{10.1109/ACCESS.2017.DOI}

\title{FPGA-Based Digital FIR Filter for Biomedical ECG Signal Denoising}
\author{\uppercase{Asrafuzzaman Khan Nahin}\authorrefmark{1}}
\address[1]{Department of Electrical Engineering, King Fahd University of Petroleum and Minerals, Dhahran 31261, Saudi Arabia (e-mail: g202318990@kfupm.edu.sa)}
\tfootnote{This work was carried out as the EE 541 course project.}

\markboth
{Khan Nahin: FPGA-Based Digital FIR Filter for ECG Denoising}
{Khan Nahin: FPGA-Based Digital FIR Filter for ECG Denoising}

\corresp{Corresponding author: Asrafuzzaman Khan Nahin (e-mail: g202318990@kfupm.edu.sa).}

\begin{abstract}
Electrocardiogram (ECG) recordings are often corrupted by high-frequency interference such as mains pickup and muscle noise. This paper presents the design, verification and FPGA implementation of a 16-bit fixed-point low-pass FIR filter that removes such noise from a prerecorded ECG segment of 1024 samples at 360~Hz. A 63-tap Hamming-window filter with a 20~Hz cutoff is designed in Python, its coefficients are quantized to signed Q15, and a bit-exact integer model with a 40-bit accumulator, one arithmetic shift and saturation serves as the reference. The Verilog implementation reuses a single multiplier over 63 clock cycles in a three-stage pipeline and needs 66 clock cycles per sample. Samples are read from block RAM and released by a 360~Hz strobe. A self-checking testbench finds zero mismatches against the Python reference for the ECG data and seven stress vectors. On a Nexys~3 board (Spartan-6) the complete design uses \resLUT{} LUTs, \resFF{} flip-flops, \resDSP{} DSP block and \resBRAM{} block RAMs, and meets 100~MHz timing with a worst negative slack of \resWNS{}~ns. Every filtered sample is sent to a PC over Ethernet and compared with the reference. The filter raises the signal-to-noise ratio from $-0.92$~dB to $6.76$~dB and lowers the noise power by 16.3~dB; the remaining error is dominated by distortion of the QRS complex caused by the 20~Hz cutoff. The paper also reports a synthesis-dependent error that appeared only on the hardware, and how it was found and removed.
\end{abstract}

\begin{keywords}
ECG denoising, FIR filter, fixed-point arithmetic, FPGA, Verilog, Xilinx Spartan-6
\end{keywords}

\titlepgskip=-15pt

\maketitle

\section{Introduction}
\label{sec:introduction}
\PARstart{T}{he} electrocardiogram (ECG) is recorded at the skin surface and has an amplitude of about one millivolt, so it is easily disturbed by power-line interference, muscle activity and electronic noise. Most of the diagnostic energy of an ECG lies below about 40~Hz, while much of this interference lies above it, which makes a low-pass finite impulse response (FIR) filter a natural first denoising stage. FIR filters are always stable, can have exactly linear phase, and map directly onto the multiply-accumulate (MAC) resources of an FPGA.

The objective of this project is to design, verify and implement a fixed-point FIR low-pass filter on an FPGA and to demonstrate that the hardware produces the same result as a software reference. The contributions are:
\begin{itemize}
\item a complete design flow from the filter specification to a board demonstration, with a bit-exact Python model of the hardware arithmetic;
\item a small, readable Verilog design that uses one multiplier over 63 clock cycles, with a self-checking testbench that includes Ethernet frame checks;
\item measured implementation results (resources, timing) and a comparison of the signal captured from the board with the reference;
\item a description of a synthesis-dependent error that the RTL simulation did not show and how it was located with a post-synthesis simulation.
\end{itemize}

Section~\ref{sec:data} describes the data, Section~\ref{sec:design} the filter design, and Section~\ref{sec:fixed} the fixed-point implementation. Section~\ref{sec:rtl} presents the RTL architecture, Section~\ref{sec:verification} the verification, Section~\ref{sec:fpga} the FPGA implementation and Section~\ref{sec:demo} the hardware demonstration. Section~\ref{sec:discussion} discusses the results.

\section{Data and Signal Model}
\label{sec:data}
The input is the prerecorded ECG segment supplied with the project, stored in \texttt{ecg\_fir\_data\_full.csv}. It has 1024 rows with sample index, time, a clean reference in millivolts and in integer units, and a noisy integer input. The time stamps give a sampling rate of $(1024-1)/2.841667 \approx 360.00$~Hz, so $f_s = 360$~Hz is used throughout. The noisy input ranges from $-7525$ to $16384$, which fits in a signed 16-bit word. The clean integer column serves as the quality reference for all error measures.

The difference between the noisy and the clean column has a standard deviation of 2329 units, larger than that of the clean signal (2067 units), so the input signal-to-noise ratio is below 0~dB. A spectrum of this difference shows that 94\,\% of the noise power lies between 40 and 180~Hz with strong components near 60~Hz, whereas only 3.4\,\% lies below 20~Hz. The clean ECG, in contrast, has 79\,\% of its power below 20~Hz. A low-pass filter with a 20~Hz cutoff is therefore well matched to this noise, which is what the project specifies.

\Figure[t!](topskip=0pt, botskip=0pt, midskip=0pt)[width=0.98\columnwidth]{ecg_signals.png}
{Clean and noisy ECG segment (1024 samples, 360~Hz).\label{fig:signals}}

\section{Filter Design}
\label{sec:design}
\subsection{Specification}
The specification is a low-pass FIR filter with a cutoff of $f_c = 20$~Hz at $f_s = 360$~Hz. The window method \cite{b1} was chosen with a Hamming window \cite{b2}. The ideal low-pass impulse response for a cutoff $\omega_c = 2\pi f_c/f_s = \pi/9$~rad is
\begin{equation}
h_d[n] = \frac{\sin\!\big(\omega_c (n-M)\big)}{\pi (n-M)}, \qquad M = \frac{N-1}{2},
\end{equation}
which is infinitely long and non-causal. It is made finite and causal by multiplying with a Hamming window
\begin{equation}
w[n] = 0.54 - 0.46\cos\!\left(\frac{2\pi n}{N-1}\right), \quad 0 \le n \le N-1,
\end{equation}
so that $h[n] = h_d[n]\,w[n]$. For an odd length $N$ and symmetric coefficients $h[n]=h[N-1-n]$ the filter has exactly linear phase and a constant group delay of $M$ samples.

\subsection{Choice of the order}
The width of the transition band of a Hamming-windowed filter is approximately $\Delta f \approx 3.3\,f_s/N$ \cite{b1}. The project does not fix the order, so $N=63$ (order 62) was chosen as a compromise between selectivity and cost. It gives $\Delta f \approx 3.3\cdot 360/63 \approx 19$~Hz, so the passband ends near 10~Hz and the stopband starts near 30~Hz around the 20~Hz cutoff. The odd length gives an integer group delay of $M=31$ samples (86~ms), which has to be compensated when comparing with the clean signal. Table~\ref{tab:spec} summarizes the design.

\begin{table}[!t]
\caption{Filter specification and measured response (floating point)}
\label{tab:spec}
\centering
\begin{tabular}{ll}
\toprule
Parameter & Value \\
\midrule
Type / window & Low-pass FIR, Hamming \\
Sampling rate $f_s$ & 360~Hz \\
Cutoff $f_c$ ($-6$~dB point) & 20~Hz (measured $-6.00$~dB) \\
Taps $N$ / order & 63 / 62 \\
Group delay & 31 samples (86.1~ms) \\
Passband ripple (0--10~Hz) & 0.034~dB \\
Stopband attenuation ($\ge$30~Hz) & 53.7~dB \\
DC gain & 0~dB (sum of coefficients $=1$) \\
\bottomrule
\end{tabular}
\end{table}

\Figure[t!](topskip=0pt, botskip=0pt, midskip=0pt)[width=0.98\columnwidth]{frequency_response.png}
{Magnitude response of the 63-tap filter. The dashed line marks the 20~Hz cutoff.\label{fig:mag}}

\Figure[t!](topskip=0pt, botskip=0pt, midskip=0pt)[width=0.98\columnwidth]{coefficients_phase.png}
{Filter coefficients (top) and unwrapped phase response (bottom). The coefficients are symmetric, so the phase is linear.\label{fig:phase}}

Figures~\ref{fig:mag} and~\ref{fig:phase} show the magnitude and phase responses. The coefficients were computed with \texttt{scipy.signal.firwin} \cite{b5}.

\section{Fixed-Point Implementation}
\label{sec:fixed}
\subsection{Number formats}
Input and output samples are signed 16-bit integers. Each coefficient is quantized to signed 16-bit Q15, that is $h_q[k] = \mathrm{round}(h[k]\cdot 2^{15})$, with values between $-527$ and $3643$. The output is
\begin{equation}
y[n] = \mathrm{sat}_{16}\!\left(\left\lfloor \frac{1}{2^{15}} \sum_{k=0}^{62} h_q[k]\,x[n-k] \right\rfloor\right),
\end{equation}
where $\lfloor\cdot\rfloor$ is an arithmetic right shift (rounding toward $-\infty$), applied once to the full sum, and $\mathrm{sat}_{16}$ limits the result to $[-32768, 32767]$. Shifting each product individually would give a different result and is avoided. Note that for negative values an arithmetic shift is not the same as truncation toward zero: $-1 \gg 15 = -1$, but $\mathrm{int}(-1/32768) = 0$. The hardware uses the arithmetic shift, so the Python model does as well.

\subsection{Coefficient quantization}
Rounding to Q15 changes each coefficient by at most half of one least significant bit, $2^{-16} = 1.53\cdot10^{-5}$. The measured maximum error is $1.45\cdot10^{-5}$ (0.48 LSB) and the root-mean-square error is $8.1\cdot10^{-6}$. Table~\ref{tab:quant} compares the floating-point and quantized responses. The two responses are practically identical: the largest difference between them is $1.6\cdot10^{-4}$ ($-76$~dB), the passband ripple is unchanged, and the stopband attenuation changes by 0.09~dB. The DC gain falls from 0 to $-0.0008$~dB because the quantized coefficients sum to 0.99991 instead of 1. Figure~\ref{fig:cmp} overlays the two magnitude responses.

\begin{table}[!t]
\caption{Effect of coefficient quantization (Q15)}
\label{tab:quant}
\centering
\begin{tabular}{lcc}
\toprule
Quantity & Float & Q15 \\
\midrule
Sum of coefficients & 1.00000 & 0.99991 \\
DC gain (dB) & 0.0000 & $-0.0008$ \\
Passband ripple, 0--10~Hz (dB) & 0.0335 & 0.0334 \\
Stopband attenuation, $\ge$30~Hz (dB) & 53.71 & 53.79 \\
Gain at 20~Hz (dB) & $-5.9971$ & $-5.9970$ \\
Max.\ coefficient error (LSB) & -- & 0.48 \\
Max.\ $|H_{float}-H_{Q15}|$ & \multicolumn{2}{c}{$1.58\cdot10^{-4}$ ($-76.0$~dB)} \\
\bottomrule
\end{tabular}
\end{table}

\Figure[t!](topskip=0pt, botskip=0pt, midskip=0pt)[width=0.98\columnwidth]{frequency_response_comparison.png}
{Magnitude response before and after quantization of the coefficients to 16 bits. The curves coincide.\label{fig:cmp}}

\subsection{Accumulator width}
The largest possible magnitude of the full sum is $X_{max}\sum_k |h_q[k]|$. With $\sum_k|h_q[k]| = 44633$ this is $16384\cdot44633 \approx 7.3\cdot10^{8}$ for the supplied data and $32768\cdot 44633 \approx 1.46\cdot10^{9}$ for any legal 16-bit input. Both are far below $2^{39}$, so the 40-bit signed accumulator cannot overflow; 32 bits would already be enough. The largest sum actually reached on the ECG data is $3.35\cdot10^{8}$. The generous width costs nothing, because a Spartan-6 DSP block has a 48-bit accumulator.

\subsection{Bit-exact reference and test vectors}
A Python model reproduces the hardware arithmetic with Python integers: a causal convolution with an initially empty history, one shift of the full sum, and saturation. It is checked against an independent NumPy integer convolution. The model produces the reference output for the ECG data and for seven additional vectors (Table~\ref{tab:vectors}), all written as 4-digit hexadecimal \texttt{.mem} files with negative values in two's complement. The fixed-point output differs from the floating-point filter by at most 1.44 LSB (0.61 LSB rms).

\section{RTL Architecture}
\label{sec:rtl}
\subsection{System overview}
Figure~\ref{fig:system} shows the complete design. A sample strobe releases one sample every $1/360$~s. The control logic in \texttt{fir\_top} reads the next noisy sample from a ROM, offers it to the FIR core and waits for the result. The result is compared with the Python reference value, stored in a second ROM, and is also sent to the PC as a UDP packet. The LEDs show the verdict.

\Figure[t!](topskip=0pt, botskip=0pt, midskip=0pt)[width=0.98\columnwidth]{diagram_system.png}
{Block diagram of the system implemented on the FPGA.\label{fig:system}}

\subsection{FIR core}
The FIR core (\texttt{fir\_mac}) uses the direct form with a sequential multiply-accumulate. Because a new sample arrives only every 277\,778 clock cycles of the 100~MHz clock, one multiplier can be reused for all 63 taps. The core holds a 63-word history of 16-bit samples, a 63-word coefficient ROM, a 16$\times$16-bit multiplier and a 40-bit accumulator. Its interface is listed in Table~\ref{tab:if}. A sample is accepted when \texttt{in\_valid} is high while the core is not busy.

\begin{table}[!t]
\caption{Interface of the FIR core}
\label{tab:if}
\centering
\begin{tabular}{lll}
\toprule
Signal & Dir. & Meaning \\
\midrule
\texttt{clk}, \texttt{rst} & in & 100~MHz clock, synchronous active-high reset \\
\texttt{in\_valid}, \texttt{x} & in & One-clock offer of a signed 16-bit sample \\
\texttt{busy} & out & High while a sample is being processed \\
\texttt{out\_valid}, \texttt{y} & out & One-clock pulse with the signed 16-bit result \\
\bottomrule
\end{tabular}
\end{table}

A four-state machine controls the core (Fig.~\ref{fig:fsm}):
\begin{itemize}
\item \textbf{IDLE}: wait for an accepted sample. The history is shifted by one place, the sample is stored at \texttt{hist[0]}, and the accumulator and tap counter are cleared. This is the hardware version of inserting the newest sample at the front of a list.
\item \textbf{MAC}: one tap is issued per clock cycle, from tap 0 (newest sample) to tap 62.
\item \textbf{FLUSH}: two cycles that let the last tap travel through the pipeline into the accumulator.
\item \textbf{OUT}: the full sum is shifted right by 15 bits with an arithmetic shift, saturated to 16 bits, and \texttt{out\_valid} is pulsed for one clock.
\end{itemize}

\Figure[t!](topskip=0pt, botskip=0pt, midskip=0pt)[width=0.98\columnwidth]{diagram_fsm.png}
{State machine of the FIR core.\label{fig:fsm}}

\subsection{Pipeline}
Doing the 63-to-1 history multiplexer, the multiplication and the 40-bit addition in one clock cycle was too slow: the first implementation had a minimum clock period of 10.378~ns and failed the 10~ns constraint. The MAC was therefore split into three stages (Fig.~\ref{fig:pipe}): the operands $h_r, c_r$ are registered, then the product $p_r$, then the sum is accumulated in \texttt{acc}. This adds three clock cycles, so one sample needs 66 cycles from acceptance to the output pulse, which is $0.66~\mu$s. This is 0.024\,\% of the 2.78~ms sample period. The pipeline registers are fed so that the product register is zero whenever no real tap is flowing; the reason is explained in Section~\ref{sec:discussion}.

\Figure[t!](topskip=0pt, botskip=0pt, midskip=0pt)[width=0.98\columnwidth]{diagram_pipeline.png}
{Three-stage MAC pipeline of the FIR core: operand registers, product register and accumulator.\label{fig:pipe}}

\subsection{Sampling strobe, memories and top level}
The \texttt{sample\_strobe} module counts 277\,778 clock cycles and then gives a one-clock pulse; this is a clock enable and not a divided clock, so the whole design stays in one 100~MHz clock domain. The ECG samples and the reference outputs are stored in two $1024\times16$ ROMs (\texttt{ecg\_rom}), which are initialized from the \texttt{.mem} files with \texttt{\$readmemh} and map to block RAM. The coefficients are loaded the same way into a small distributed ROM.

\texttt{fir\_top} runs the sequence strobe $\rightarrow$ read sample $\rightarrow$ filter $\rightarrow$ compare $\rightarrow$ send for each of the 1024 samples, counts the mismatches and shows the result on the LEDs (LD2 = pass, LD3 = fail, LD4--LD7 = mismatch count). After a pause of 0.67~s the run repeats and the filter history is cleared, so that a PC can join at any time. Each result is sent as one UDP broadcast packet (20-byte payload: marker, flags, sample number, filtered value, reference value, noisy input) by a transmitter taken from an earlier project that was verified on the same board.

\section{Verification}
\label{sec:verification}
\subsection{Self-checking testbenches}
Two testbenches were written. \texttt{tb\_fir\_mac} drives the filter core with every sample of an input file, waits for each result, compares it with the corresponding word of the Python reference, counts mismatches and prints PASS or FAIL. It also checks that there is exactly one output pulse per sample and that \texttt{busy} is never unknown. \texttt{tb\_fir\_top} tests the complete design: it additionally plays the role of the Ethernet PHY, decodes every transmitted frame, and checks the Ethernet type, the UDP port, the CRC, the sample number and the values; it verifies that the run restarts by itself and that a second run produces identical output.

Seven stress vectors, in addition to the ECG data, are used to test cases that the ECG data never reaches (Table~\ref{tab:vectors}). In particular, the two clamp vectors use input signs that follow the signs of the coefficients, which is the worst case of Section~\ref{sec:fixed}, so the output is driven beyond the 16-bit range and saturation is exercised (7 clamped samples each).

\begin{table}[!t]
\caption{Test vectors (1024 samples each) and RTL results}
\label{tab:vectors}
\centering
\footnotesize
\begin{tabular}{llc}
\toprule
Vector & Purpose & Mismatches \\
\midrule
ECG (noisy input) & Real data & 0 \\
Impulse (32767, then 0) & Replays the coefficients & 0 \\
Positive step (+10000) & Settles at 9999 & 0 \\
Negative step ($-10000$) & Settles at $-10000$ & 0 \\
Extreme ($\pm$full scale, alternating) & Worst-case fast input & 0 \\
Random (full range, seeded) & General check & 0 \\
Clamp positive & Saturation at $+32767$ & 0 \\
Clamp negative & Saturation at $-32768$ & 0 \\
\bottomrule
\end{tabular}
\end{table}

The RTL matches the Python reference exactly on all eight inputs. To show that the testbench can fail, deliberate faults were introduced: changing the shift from 15 to 14 gave 1023 mismatches, a lowered clamp limit was detected, and wrong reference data was detected by the full-system testbench (15 or more mismatches shown on the LEDs, FAIL lit).

\subsection{Waveforms}
Figure~\ref{fig:wave} shows a simulation of one sample. The \texttt{in\_valid} pulse is followed by the \texttt{MAC} state during which \texttt{tap\_count} runs from 0 to 62, then \texttt{FLUSH} and \texttt{OUT}; \texttt{out\_valid} pulses once and \texttt{y} takes its new value. \texttt{busy} is high from acceptance until the output pulse.

\Figure[t!](topskip=0pt, botskip=0pt, midskip=0pt)[width=0.98\columnwidth]{sim_waveform.png}
{Simulation waveform of the FIR core for one sample (clock, input handshake, state, tap counter, output handshake and output value).\label{fig:wave}}

\subsection{Denoising performance}
The quality of the filtering is measured against the clean ECG. Because the filter delays the signal by 31 samples, the output $y[n]$ is compared with the clean value $c[n-31]$. The mean squared error (MSE) and the signal-to-noise ratio (SNR) are
\begin{equation}
\mathrm{MSE}=\frac{1}{L}\sum_n \big(c[n-31]-\hat{s}[n]\big)^2,\quad
\mathrm{SNR}=10\log_{10}\frac{\sum c^2}{\sum (c-\hat{s})^2},
\end{equation}
where $\hat s$ is either the noisy input or the filter output over $L=993$ samples. Table~\ref{tab:metrics} lists the results of the bit-exact model and of the output captured from the FPGA (Section~\ref{sec:demo}); they are identical. The MSE falls by a factor of 5.9 and the SNR improves by 7.7~dB.

\begin{table}[!t]
\caption{Denoising results against the clean ECG (integer ADC units)}
\label{tab:metrics}
\centering
\begin{tabular}{lcc}
\toprule
Signal & MSE & SNR (dB) \\
\midrule
Noisy input & $5.42\cdot10^{6}$ & $-0.92$ \\
Filtered, fixed-point model (Python) & $9.26\cdot10^{5}$ & $6.76$ \\
Filtered, captured from the FPGA & $9.26\cdot10^{5}$ & $6.76$ \\
\bottomrule
\end{tabular}
\end{table}

\Figure[t!](topskip=0pt, botskip=0pt, midskip=0pt)[width=0.98\columnwidth]{filtered_ecg_signal.png}
{Clean ECG and filtered ECG (floating-point filter, delay removed).\label{fig:filtered}}

\section{FPGA Implementation}
\label{sec:fpga}
The target board is a Digilent Nexys~3 with a Xilinx Spartan-6 XC6SLX16 \cite{b3}. The design was synthesized, placed and routed with Xilinx ISE 14.6 using the constraints file \texttt{fir.ucf}, which fixes the 100~MHz clock, the button, switches, LEDs and the Ethernet PHY pins. ISE maps the multiplier onto one DSP48A1 slice \cite{b4} and the ROMs onto block RAM. Table~\ref{tab:res} lists the resources and the timing reported by the tools for the complete design including the Ethernet transmitter.

\begin{table}[!t]
\caption{Implementation results on XC6SLX16 (100~MHz constraint)}
\label{tab:res}
\centering
\begin{tabular}{lcc}
\toprule
Resource & Used & Available \\
\midrule
Slice LUTs & \resLUT{} (13\,\%) & 9112 \\
Slice registers (flip-flops) & \resFF{} (10\,\%) & 18224 \\
DSP48A1 & \resDSP{} (3\,\%) & 32 \\
Block RAM (RAMB16) & \resBRAM{} (6\,\%) & 32 \\
Bonded IOBs & \resIOB{} (9\,\%) & 232 \\
\midrule
Minimum clock period & \multicolumn{2}{c}{\resPeriod{}~ns ($f_{max}\approx\resFmax$~MHz)} \\
Worst negative slack (10~ns) & \multicolumn{2}{c}{\resWNS{}~ns} \\
Timing errors & \multicolumn{2}{c}{0 (all constraints met)} \\
\bottomrule
\end{tabular}
\end{table}

The timing constraints are met with a positive worst slack. The first version, in which the history multiplexer, multiplier and adder were in one clock cycle, had a minimum period of 10.378~ns and violated the constraint; the pipeline of Section~\ref{sec:rtl} removed the violation. Power consumption was not measured.

\section{Hardware Demonstration}
\label{sec:demo}
The bitstream was loaded into the board over JTAG. The Nexys~3 has a USB-UART bridge, but its serial port did not appear on the host computer, so the filtered ECG is returned over the Ethernet port instead, which was verified with an earlier design on the same board. The PC capture program \texttt{fir\_capture.py} listens for the UDP packets, stores the 1024 samples in a CSV file, compares each value with the Python reference and computes the MSE and SNR against the clean ECG.

The measurement was repeated three times, starting the capture at different points of the repeating run. In every capture:
\begin{itemize}
\item all 1024 different sample numbers were received (no packet lost, at 360 packets per second);
\item the noisy input and the reference value carried in the packets equal \texttt{ecg.mem} and \texttt{expected.mem};
\item \textbf{the filtered value from the FPGA equals the Python reference at all 1024 samples (0 differences).}
\end{itemize}
Computing the MSE and SNR from the captured data against the clean ECG gives exactly the values of the Python model (Table~\ref{tab:metrics}): MSE $9.26\cdot10^{5}$ and SNR $6.76$~dB, compared with $5.42\cdot10^{6}$ and $-0.92$~dB for the noisy input. Figure~\ref{fig:capture} shows the captured signal. The high-frequency noise is removed and the baseline follows the clean ECG, while the R peaks are lower than in the clean signal because of the 20~Hz cutoff (Section~\ref{sec:discussion}). The output is delayed by 31 samples, which is removed in the figure.

The design also compares every output with the reference on the board and shows the verdict on the LEDs (LD2 = pass); in this work the outcome was confirmed from the captured data. At 360 samples per second one run of 1024 samples takes 2.84~s, and the run repeats after a pause of 0.67~s.

\Figure[t!](topskip=0pt, botskip=0pt, midskip=0pt)[width=0.98\columnwidth]{fpga_capture.png}
{Filtered ECG captured from the FPGA over Ethernet (red) compared with the clean ECG (dashed) and the noisy input (grey). The output is shifted back by the 31-sample filter delay.\label{fig:capture}}

\section{Discussion}
\label{sec:discussion}
\subsection{Why the SNR gain is limited}
The filter reduces the noise power from $5.42\cdot10^{6}$ to $1.26\cdot10^{5}$, which is 16.3~dB, but the SNR improves by only 7.7~dB. The reason is that the 20~Hz cutoff also removes part of the ECG itself: 21\,\% of the clean signal power lies above 20~Hz, mainly in the sharp QRS complex, and the filter keeps only 69\,\% of the signal power. The error after filtering consists of $8.0\cdot10^{5}$ from this signal distortion and only $1.3\cdot10^{5}$ from the remaining noise. Denoising and preserving the waveform are therefore conflicting goals. An additional Python experiment with the same 63-tap Hamming design shows how the result depends on the cutoff (Table~\ref{tab:sweep}); the project specifies 20~Hz, but a cutoff of 35--40~Hz would give about 11.4~dB for this noise, which lies mostly above 40~Hz.

\begin{table}[!t]
\caption{SNR versus cutoff for the same 63-tap Hamming design (Python, floating point)}
\label{tab:sweep}
\centering
\begin{tabular}{lcccccc}
\toprule
Cutoff (Hz) & 15 & 20 & 25 & 30 & 35 & 40 \\
\midrule
SNR (dB) & 4.57 & 6.76 & 8.88 & 10.50 & 11.32 & 11.48 \\
\bottomrule
\end{tabular}
\end{table}

\subsection{A synthesis-dependent error found on the board}
The first complete design passed all simulations, but the data captured from the board differed from the reference at 861 of 1024 samples. The differences were small (up to 13 LSB) after the first 62 samples. Fitting the captured output as a linear function of the input history showed that the board behaved like the intended 63-tap filter plus one extra 64th tap with the coefficient of the last tap acting on the oldest sample. In other words, one stale product from the previous sample was added to every sum. The reason was found with a simulation of the netlist produced by the synthesis tool: ISE had merged the multiplier, the product register and the accumulator into one DSP block, and had implemented ``clear the accumulator when a sample is accepted'' as ``replace the old sum by the product register'' instead of by zero. In the RTL simulation the product register is never used at that moment, so the simulation could not show the problem. The post-synthesis simulation reproduced the 800 mismatches of the hardware, and the correction was tested in the same way (0 mismatches) before building the bitstream again. The correction forces the product register to zero whenever no real tap is flowing, so the result no longer depends on how the clear is implemented.

A second difference was seen in the first 62 samples of each repeated run: the filter history still held the end of the previous run, whereas the Python reference starts from an empty history. The full-system testbench was extended to check a second run (it failed with 62 mismatches before the correction), and the top level now clears the filter at the start of every run.

These two findings show why a simulation of the RTL alone is not sufficient evidence: the final check has to compare the signal that really leaves the board with the reference.

\subsection{Limitations}
The design processes a stored ECG segment and not a live ADC signal; a live input would only replace the ROM by an input register. The DSP block, the ROMs and the Ethernet transmitter are inferred or reused, and power consumption and the effect of long operation were not measured. The serial (UART) route to the PC was attempted but the board's UART port was not enumerated by the host; the Ethernet route was used instead. The filter order and window were chosen by the author and not optimized; an equiripple design could meet the same specification with fewer taps.

\section{Conclusion}
A 63-tap, 16-bit fixed-point low-pass FIR filter for ECG denoising was designed, quantized, implemented in Verilog and run on a Nexys~3 FPGA. The RTL matches a bit-exact Python model on the ECG data and seven stress vectors; the design meets 100~MHz timing with \resWNS{}~ns slack using \resLUT{} LUTs, \resFF{} flip-flops, \resDSP{} DSP block and \resBRAM{} block RAMs; and it needs 66 clock cycles per sample against 277\,778 available. The filter lowers the noise power by 16.3~dB and improves the SNR by 7.7~dB. The study also showed that matching the reference in RTL simulation is not enough: a synthesis-dependent error was found only by capturing the output of the board, and was located with a post-synthesis simulation. All sources, scripts and the build instructions are in the project archive.

\begin{thebibliography}{00}
\bibitem{b1} A. V. Oppenheim and R. W. Schafer, \emph{Discrete-Time Signal Processing}, 3rd ed. Upper Saddle River, NJ, USA: Pearson, 2010.
\bibitem{b2} F. J. Harris, ``On the use of windows for harmonic analysis with the discrete Fourier transform,'' \emph{Proc. IEEE}, vol. 66, no. 1, pp. 51--83, Jan. 1978.
\bibitem{b3} Digilent Inc., ``Nexys3 FPGA Board Reference Manual,'' Doc.\ 502-182.
\bibitem{b4} Xilinx Inc., ``Spartan-6 FPGA DSP48A1 Slice User Guide,'' UG389.
\bibitem{b5} P. Virtanen \emph{et al.}, ``SciPy 1.0: Fundamental algorithms for scientific computing in Python,'' \emph{Nature Methods}, vol. 17, pp. 261--272, 2020.
\end{thebibliography}

\EOD
\end{document}
